\documentclass[12pt,a4paper]{article}
\usepackage[utf8]{inputenc}
\usepackage[T1]{fontenc}
\usepackage[magyar]{babel}
\usepackage{lmodern}
\usepackage{geometry}
\usepackage{hyperref}
\usepackage{xcolor}
\usepackage{listings}
\usepackage{setspace}

\geometry{margin=2.5cm}
\onehalfspacing
\setlength{\parskip}{0.6em}
\setlength{\parindent}{0pt}

\hypersetup{
  colorlinks=true,
  linkcolor=blue,
  urlcolor=blue
}

\lstset{
  basicstyle=\ttfamily\small,
  frame=single,
  breaklines=true,
  columns=fullflexible,
  backgroundcolor=\color{black!2},
  keywordstyle=\color{blue!70!black},
  commentstyle=\color{gray!70!black}
}

\title{IT biztonság}
\author{Gyurkovics Dominik Balázs\\Debreceni Egyetem Informatikai Kar}
\date{\today}

\begin{document}
\maketitle
\tableofcontents
\newpage

\section{Bevezetés}
Az IT biztonság célja, hogy a rendszerek és adatok a teljes életciklusuk során megőrizzék a \textbf{bizalmasságot}, \textbf{sértetlenséget} és \textbf{rendelkezésre állást}. A gyakorlati védelem nem egyetlen eszközön múlik, hanem egymásra rétegzett megoldásokon: jogosultságkezelés, hitelesítés, kriptográfia, hálózati védelem, monitoring és helyreállítási eljárások. A következő fejezetek összefoglalják a Linux és Windows környezetben alkalmazott, tananyagban szereplő megoldásokat, és a parancssori példák segítségével kézzelfoghatóvá teszik a folyamatokat.

\section{Rendszer-helyreállítás és alapadminisztráció}
Vészhelyzetben a rendszerhez való hozzáférés biztosítása kritikus. Linux környezetben a \texttt{rescue} mód és a \texttt{rd.break} kernelparaméter teszi lehetővé, hogy a root fájlrendszert írhatóvá téve helyreállítsuk a jelszavakat vagy a hibás konfigurációkat. A helyreállítás során figyelni kell a SELinux címkék visszaállítására is, különben később működési problémák jelentkezhetnek.

\begin{lstlisting}[language=bash,caption={Helyreállítási műveletek (cheatsheet részlet)}]
rd.break
mount -o remount,rw /sysroot
chroot /sysroot
passwd
touch ./autorelabel
\end{lstlisting}

Ezek a lépések biztosítják, hogy a rendszer újraindításkor automatikusan újracímkézze a fájlokat, és a hozzáférések ismét konzisztens állapotba kerüljenek.

\section{Felhasználók, jelszavak és hitelesítés}
A Linux felhasználói adatok a \texttt{/etc/passwd}, \texttt{/etc/group} és \texttt{/etc/shadow} állományokban találhatók. Míg a \texttt{passwd} a felhasználó azonosítóit és beállításait tartalmazza, a tényleges jelszóhash a \texttt{shadow} fájlba kerül. A jelszó élettartama, a figyelmeztetési időszak és a fiók lejárata mind auditálható és szigorítható, ami alapkövetelmény a biztonságos üzemeltetésben.

\begin{lstlisting}[language=bash,caption={Felhasználó és jelszó-életciklus kezelés}]
useradd user1
id user1
passwd -l user1
chage -l user1
chage -d 2020-01-01 -m 0 -M 30 -I 30 -E 2020-04-01 user1
\end{lstlisting}

\subsection{Lejárati dátum és jelszó-élettartam finomhangolás}
A fiók lejáratát és a jelszavak cseréjének kényszerítését a \texttt{chage} parancs kezeli. A dátumformátum \texttt{YYYY-MM-DD}, és érdemes figyelmeztetési időszakot (\texttt{-W}) és inaktív időszakot (\texttt{-I}) is beállítani, hogy a felhasználó még a fiók lezárása előtt értesítést kapjon.

\begin{lstlisting}[language=bash,caption={Lejárati dátum és figyelmeztetések}]
chage -E 2026-12-31 user1
chage -W 7 -I 14 user1
passwd --expire user1
chage -l user1
\end{lstlisting}

\subsection{SSH és kulcsalapú autentikáció}
Az SSH protokoll titkosított csatornát biztosít a távoli adminisztrációhoz. A jelszavas autentikáció helyett az aszimmetrikus kulcspár használata javasolt, ahol a privát kulcs kizárólag a felhasználónál, a publikus kulcs pedig a szerveren található. A kulcspár használata egyszerre növeli a biztonságot és egyszerűsíti a hozzáférés menedzsmentjét.

\begin{lstlisting}[language=bash,caption={Kulcspár létrehozás és telepítés}]
ssh-keygen -t rsa -b 4096
ssh-copy-id ssh-user-3@localhost
chmod 700 ~/.ssh
chmod 600 ~/.ssh/authorized_keys
\end{lstlisting}

\section{Hozzáférés-szabályozás és jogosultságok}
Linux rendszereken az alapvető hozzáférés-szabályozás a klasszikus DAC (Discretionary Access Control) modellen alapul. Az \texttt{rwx} jogosultságok az \textit{owner}, \textit{group} és \textit{other} körök szerint értelmezhetők, katalógusok esetében pedig eltérő jelentésük van (pl. \textit{execute} nélkül a katalógus használhatatlan). A jogosultságok finomhangolására az ACL (Access Control List) szolgál, amely további felhasználók és csoportok számára is képes jogokat rendelni.

\begin{lstlisting}[language=bash,caption={Jogosultságok és ACL kezelése}]
ls -l
chmod 640 file.txt
chown user:group file.txt
getfacl file.txt
setfacl -m u:user:rwx file.txt
\end{lstlisting}

\subsection{Numerikus chmod értékek értelmezése}
A numerikus jelölésnél minden számjegy három bitet kódol: olvasás (4), írás (2), végrehajtás (1). A három számjegy rendre a tulajdonos, csoport és más felhasználók jogait jelöli. Négy számjegy esetén az első a speciális biteket tartalmazza: \texttt{4=setuid}, \texttt{2=setgid}, \texttt{1=sticky}.

\begin{lstlisting}[language=bash,caption={chmod gyors táblázat}]
chmod 000 file.txt  # --- --- ---
chmod 640 file.txt  # rw- r-- ---
chmod 750 dir/      # rwx r-x ---
chmod 1777 /tmp     # sticky bit, mindenki írhat, törlés védett
\end{lstlisting}

A speciális bitek (setuid, setgid, sticky) további biztonsági kontrollt adnak: a \textit{sticky bit} például megakadályozza, hogy bárki törölni tudjon fájlokat egy megosztott katalógusból. A rendszerszintű kontrollt a SELinux biztosítja, amely MAC (Mandatory Access Control) elven működik, és akár az ACL által engedélyezett műveleteket is megtagadhatja.

Windows környezetben a NTFS jogosultságok és az \textit{Effective Access} vizsgálata segít megérteni, hogy egy adott felhasználó mit érhet el, míg az Auditing naplózza a hozzáférési eseményeket. Ez az ellenőrizhetőség elengedhetetlen a megfelelőség (compliance) és az incidensek kivizsgálása során.

\section{Fájlrendszer létrehozása és partíciókezelés}
A biztonságos tárolás alapja a kontrollált partícionálás és csatolás. A következő lépések bemutatják, hogyan készítsünk új fájlrendszert a \texttt{/dev/sdc} lemezen, hogyan hozzuk létre a \texttt{/dev/sdc1} partíciót, és hogyan csatoljuk azt egy fix mount pontra. A \texttt{blkid} parancs segítségével a tartós csatoláshoz szükséges UUID is megadható.

\begin{lstlisting}[language=bash,caption={Partíció létrehozás /dev/sdc-n}]
lsblk -f
sudo wipefs -a /dev/sdc
sudo parted -s /dev/sdc mklabel gpt
sudo parted -s /dev/sdc mkpart primary ext4 1MiB 100%
sudo partprobe /dev/sdc
lsblk /dev/sdc
sudo mkfs.ext4 -L data /dev/sdc1
sudo mkdir -p /mnt/sdc1
sudo mount /dev/sdc1 /mnt/sdc1
sudo blkid /dev/sdc1
sudo sh -c 'echo "UUID=... /mnt/sdc1 ext4 defaults,noatime 0 2" >> /etc/fstab'
\end{lstlisting}

\section{Hálózatbiztonság és tűzfalak}
A tűzfalak elsődleges szerepe a bejövő és kimenő hálózati forgalom szűrése. A host-based tűzfalak (pl. Linux Netfilter, Windows Firewall) közvetlenül az operációs rendszer szintjén védenek, és a csomagok fejlécadatai, illetve a kapcsolati állapot (stateful inspection) alapján hoznak döntést. A Linux \texttt{firewalld} zónákon keresztül egyszerűsíti a szabályok kezelését, míg Windowsban profilok (Domain, Private, Public) szerint történik a konfiguráció.

\begin{lstlisting}[language=bash,caption={Firewalld alapműveletek}]
ss -ltnp
firewall-cmd --state
firewall-cmd --zone=internal --add-service=http
firewall-cmd --runtime-to-permanent
nft list ruleset
\end{lstlisting}

A hálózati forgalom megfigyelésére a rögzítés és elemzés ad lehetőséget. A \texttt{dumpcap} vagy a Wireshark a forgalom mintavételezésével segít észlelni a gyanús aktivitást, különösen incidenskezelési helyzetekben.

\begin{lstlisting}[language=bash,caption={Forgalom rögzítése és szűrése}]
dumpcap -i enp0s3 -w /tmp/1.pcap -a duration:30
dumpcap -i enp0s3 -f "host 10.0.0.1 && (icmp[0] == 0 || icmp[0] == 8 )"
\end{lstlisting}

\section{Kriptográfia és kulcskezelés}
A kriptográfia alapja, hogy a biztonságos adatkezeléshez megfelelő algoritmusokat és kulcsmenedzsmentet alkalmazzunk. Szimmetrikus titkosításnál (pl. AES) ugyanaz a kulcs szükséges a titkosításhoz és visszafejtéshez. Az ECB mód egyszerű, de mintázatot hagyhat az adatokban, ezért a CBC mód, só és inicializáló vektor használatával lényegesen biztonságosabb megoldást nyújt.

\begin{lstlisting}[language=bash,caption={OpenSSL AES és RSA parancsok}]
openssl aes-128-cbc -k PASSWORD -in file -out file.aes -p
openssl aes-128-cbc -d -k PASSWORD -in file.aes
openssl genpkey -algorithm RSA -out key.pem -pkeyopt rsa_keygen_bits:2048
openssl dgst -sign key.pem -sha256 -out document.sign document
openssl dgst -verify key.pem.pub -sha256 -signature document.sign document
\end{lstlisting}

\subsection{Publikus kulcsos titkosítás vágólapról}
Ha a publikus kulcs (PEM formátum, \texttt{-----BEGIN PUBLIC KEY-----}) a vágólapon van, először mentsük fájlba, majd ellenőrizzük és használjuk titkosításhoz. Nagy fájlok esetén ajánlott hibrid megoldást alkalmazni: a tartalmat szimmetrikus kulccsal titkosítani, a kulcsot pedig publikus kulccsal védeni.

\begin{lstlisting}[language=bash,caption={Publikus kulcs mentése és titkosítás}]
cat > public.pem
chmod 600 public.pem
openssl pkey -pubin -in public.pem -text -noout
openssl pkeyutl -encrypt -pubin -inkey public.pem -in secret.txt -out secret.txt.enc
openssl rand -out sym.key 32
openssl enc -aes-256-gcm -salt -pbkdf2 -iter 200000 -in secret.txt -out secret.txt.enc -pass file:./sym.key
openssl pkeyutl -encrypt -pubin -inkey public.pem -in sym.key -out sym.key.enc
\end{lstlisting}

Az aszimmetrikus eljárások (pl. RSA) megoldják a kulcsterjesztés problémáját, és lehetővé teszik a digitális aláírást. A tanúsítványok a nyilvános kulcsot hitelesítik, így garantálják a kapcsolat bizalmasságát és azonosíthatóságát. A tanúsítványlánc és a trust store biztosítja, hogy a kliens a megfelelően aláírt és érvényes kulcsokban bízzon.

\section{Titkosított tárolók és platformbiztonság}
A fájlrendszer szintű titkosítás a fizikai hozzáférés kockázatát kezeli. A Linux LUKS szabvány lehetővé teszi, hogy egy meghajtó titkosítását több jelszóval vagy kulcsfájllal oldjuk fel, miközben a mesterkulcs védetten tárolódik. A csatolási és feloldási folyamat automatizálható a \texttt{/etc/fstab} és \texttt{/etc/crypttab} állományok segítségével.

\begin{lstlisting}[language=bash,caption={LUKS beállítási példa}]
cryptsetup luksFormat /dev/sdb1
cryptsetup luksOpen /dev/sdb1 secret
mkfs.xfs /dev/mapper/secret
mount /dev/mapper/secret /home
cat /etc/crypttab
\end{lstlisting}

A TPM (Trusted Platform Module) hardveres védelmet biztosít, a PCR regiszterek segítségével azonosítja a boot folyamat integritását, és csak akkor adja ki a titkosított adatokat, ha a rendszer állapota változatlan. Windows környezetben a BitLocker hasonló elven működik: a FVEK és VMK kulcsok kombinációjával titkosít, és a TPM-re támaszkodva védi az adatokat. A helyreállítási kulcs készítése kritikus, mert hardverváltozás esetén a TPM megakadályozhatja az automatikus feloldást.

\section{Megfigyelés és auditálás}
A biztonságos üzemeltetés egyik alappillére az események nyomon követése. Linux rendszerekben a naplófájlok és audit rendszerek, Windowsban az Event Viewer biztosítja a hozzáférések és műveletek visszakövethetőségét. Az auditálás lehetővé teszi, hogy incidensek esetén pontosan rekonstruáljuk az események láncolatát, és igazolni tudjuk a szabályok betartását.

\begin{lstlisting}[language=bash,caption={Windows és Linux audit eszközök}]
gpedit.msc
eventvwr.msc
\end{lstlisting}

\section{Összegzés}
Az IT biztonság nem egyszeri feladat, hanem folyamatos kockázatkezelési és üzemeltetési tevékenység. A jogosultságok helyes beállítása, a biztonságos hitelesítés, a jól konfigurált tűzfalak, a kriptográfiai eljárások és a titkosított tárolók mind együttesen biztosítják, hogy a rendszer ellenálljon a fenyegetéseknek. A fenti eszközök és eljárások alkalmazásával olyan védelmi rétegek építhetők, amelyek a modern informatikai környezetben elvárt biztonsági szintet nyújtják.

\appendix
\section{Online források és dokumentációk}
Az alábbi linkek megbízható, naprakész dokumentációkat tartalmaznak a fenti parancsokról és eljárásokról:

\begin{itemize}
  \item \href{https://man7.org/linux/man-pages/}{Linux man oldalak (man7.org)}
  \item \href{https://www.openssh.com/manual.html}{OpenSSH kézikönyvek}
  \item \href{https://www.openssl.org/docs/manmaster/man1/}{OpenSSL parancs dokumentáció}
  \item \href{https://firewalld.org/documentation.html}{firewalld dokumentáció}
  \item \href{https://wiki.nftables.org/wiki-nftables/index.php/Main_Page}{nftables wiki}
  \item \href{https://www.freedesktop.org/software/systemd/man/}{systemd és systemctl man oldalak}
  \item \href{https://www.kernel.org/doc/html/latest/}{Linux kernel dokumentáció}
\end{itemize}

\section{CHEATSHEET (teljes, emberileg olvasható formában)}
A következő melléklet tartalmazza a frissített \texttt{CHEATSHEET.md} teljes szövegét, kiegészítve a részletes opciólistákkal és példákkal.

\lstinputlisting[language=bash]{../CHEATSHEET.md}

\end{document}
